\chapter{Теоретическая часть}\label{ch:theor_ch}
В данной главе формализуется рассматриваемая производственная задача и обосновывается выбор математической модели.
Сначала описывается технологический контекст целлюлозно-бумажного производства и формулируются ключевые требования к плану выпуска и нарезки.
Далее приводится обзор литературы по классической задаче линейного раскроя и её отраслевым обобщениям, включая постановки с переналадками, календарным планированием и складскими ограничениями.
Завершает главу постановка задачи: вводятся обозначения, критерии оптимальности и набор ограничений, определяющих выполнимость глобального плана.

\section{Описание задачи}\label{sec:theor_ch/problem_description}
В разделе описывается производственная цепочка и фиксируются особенности постановки, которые должны быть отражены в математической модели.
Для удобства восприятия сначала вводятся основные термины и элементы технологической схемы, затем формулируются критерии качества плана и перечень факторов, учитываемых в модели.

% \subsection{Технологический контекст и элементы производственной цепочки}\label{subsec:theor_ch/problem_context}
Существует потребность в производстве рулонов бумаги, например на целлюлозно-бумажных комбинатах.
Заказы по рулонам бумаги имеют множество параметров: ширину, количество рулонов, срок, марка бумаги, плотность и количество слоев.

\begin{figure}
    \resizebox{\textwidth}{!}{
        \begin{tikzpicture}[
    node distance=0.8cm and 0.8cm,
    box/.style={
        draw,
        very thick,
        minimum width=7.2cm,
        minimum height=2.8cm,
        align=center,
        font=\small
    },
    smallbox/.style={
        draw,
        rounded corners,
        align=center,
        font=\scriptsize,
        inner sep=3pt,
        fill=gray!5
    },
    arrow/.style={-Latex, semithick},
    lbl/.style={font=\scriptsize, fill=white, inner sep=1pt}
]

% Центральная функция IDEF0
\node[box] (func) at (0,0) {\textbf{A0}\\[2mm]
Сформировать и выполнить\\
производственный план ЦБК};

% Входы (слева)
\node[smallbox] (orders) at (-6,0.8) {Заказы на бумажную\\продукцию};
\node[smallbox] (raw) at (-6,-0.9) {Исходные материалы\\и полуфабрикаты};

% Управление (сверху)
\node[smallbox] (ctrl1) at (-2.8,3.0) {Сроки выполнения\\заказов};
\node[smallbox] (ctrl2) at (0.8,3.0) {Технологические\\ограничения};
\node[smallbox] (ctrl3) at (4.2,3.0) {Производственные\\нормативы};

% Выходы (справа)
\node[smallbox] (out1) at (6,0.9) {Готовая\\продукция};
\node[smallbox] (out2) at (6,-0.5) {Производственный план\\и расписание};
\node[smallbox] (out3) at (6,-1.8) {Отходы};

% Механизмы (снизу)
\node[smallbox] (mech1) at (-4.8,-3.0) {Бумагоделательные\\машины};
\node[smallbox] (mech2) at (-1.8,-3.0) {Склады\\полуфабрикатов};
\node[smallbox] (mech3) at (1.2,-3.0) {Продольно-\\резательные станки};
\node[smallbox] (mech4) at (4.4,-3.0) {Бобино-\\резательные станки};

% Стрелки: входы
\draw[arrow] (orders.east) -- node[midway, above, lbl] {} (func.west|-orders.east);
\draw[arrow] (raw.east) -- node[midway, below, lbl] {} (func.west|-raw.east);

% Стрелки: управление
\draw[arrow] (ctrl1.south) -- (func.north|-ctrl1.south);
\draw[arrow] (ctrl2.south) -- (func.north|-ctrl2.south);
\draw[arrow] (ctrl3.south) -- (func.north|-ctrl3.south);

% Стрелки: выходы
\draw[arrow] (func.east|-out1.west) -- (out1.west);
\draw[arrow] (func.east|-out2.west) -- (out2.west);
\draw[arrow] (func.east|-out3.west) -- (out3.west);

% Стрелки: механизмы
\draw[arrow] (mech1.north) -- (func.south|-mech1.north);
\draw[arrow] (mech2.north) -- (func.south|-mech2.north);
\draw[arrow] (mech3.north) -- (func.south|-mech3.north);
\draw[arrow] (mech4.north) -- (func.south|-mech4.north);

\end{tikzpicture}
    }
    \caption{Технологическая цепочка производства рулонов бумаги}
    \label{fig:production_chain}
\end{figure}

Процесс производства заказов бумаги представляет собой сложную многостадийную технологическую цепочку, включающую подготовку целлюлозной массы, формирование и обработку бумажного полотна, а также его последующую нарезку на рулоны заданных форматов \cite{hubbe2021_energy}.
Производство начинается с подачи целлюлозы и древесной массы в смесительные бассейны, где образуется однородная бумажная масса \cite{ippta_paper_and_water_2022, kocurek1983_pulp_paper_alkaline}.
Далее эта масса поступает на бумагоделательную машину (БДМ), где проходит последовательное обезвоживание, сушку и каландрирование, обеспечивая требуемую плотность и текстуру полотна.
В завершении формируется тамбур — барабанный рулон для последующей переработки.
БДМ может производить тамбуры только фиксированной длины, с заданной скоростью.
На смену марки или плотности производимых тамбуров требуется время.
Уже на этом этапе проявляется связь будущей задачи раскроя с параметрами основного производства. Ширина выпускаемого тамбура определяется характеристиками БДМ и не может свободно подстраиваться под текущий набор заказов, поэтому последующий этап нарезки всегда выполняется в условиях заранее заданной ширины полуфабриката. Кроме того, смена марки и плотности бумаги влияет не только на производительность БДМ, но и на допустимую последовательность выпуска, поскольку частые переходы между технологическими режимами увеличивают простои и потери сырья. Следовательно, задача планирования не сводится к независимому подбору схем раскроя: она должна учитывать, какой именно полуфабрикат и в какой последовательности целесообразно произвести.

На следующем этапе заготовки поступают на продольно-резательные станки (ПРС), которые нарезают тамбур на рулоны нужных форматов \cite{sgp_ccp2021_cutting_stock_modes}.
В случае, если ПРС ограничены рядом технологических параметров: максимальной шириной обрабатываемого тамбура \cite{minimizing_waste_multilevel_manufacturing_2025}, количеством доступных ножей \cite{cyberleninka_slitting_params} и временем, необходимым для их перестановки \cite{martina2022, munawar2005integration}.
При нарезке тамбуров необходимо срезать края тамбуров.
Такие края называются кромкой и они менее качественны чем бумага в центре.
Это уменьшает доступную для нарезки ширину тамбура и требует установку ножей по краям на ПРС.
С производственной точки зрения именно участок резки является местом, где требования отдельных заказов непосредственно преобразуются в конкретные схемы раскроя. Здесь одновременно проявляются ограничения по ширине, по числу ножей, по допустимому составу форматов в одном раскрое и по трудоёмкости переналадки. Поэтому даже при фиксированном полуфабрикате существует множество альтернативных способов выполнить один и тот же набор заказов, причём лучший по отходам вариант не всегда будет лучшим по времени работы участка или числу смен ножей. Это обстоятельство и делает задачу многокритериальной уже на уровне описания технологического процесса.

Некоторые заказы могут иметь больше одного слоя, тогда на ПРС подается столько же тамбуров, сколько и требуется слоев.
БДМ может производить только однослойные тамбуры.

Также существуют бобино-резательные станки (БРС), которые имеют такие же параметры как и ПРС, но нарезают не тамбуры с БДМ, а рулоны, которые были нарезаны на других ПРС.
БРС используются для нарезки заказов с небольшой шириной.
Введение БРС означает, что в производственной цепочке появляется дополнительный уровень преобразования полуфабриката. Часть заказов не может или экономически нецелесообразно выполняться непосредственной нарезкой тамбуров на ПРС, поэтому сначала формируются промежуточные рулоны, а затем уже они перерабатываются на БРС в продукцию требуемой ширины. Это существенно усложняет постановку по сравнению с классической задачей линейного раскроя, где предполагается единственный этап разрезания исходного материала. В рассматриваемом случае план должен учитывать, какие заказы целесообразно направить на ПРС, какие --- на БРС, и как связать между собой оба этапа по времени и по доступности промежуточных рулонов.

Если все ПРС заняты или требуется накопить несколько тамбуров для нарезки заказа, тамбуры и рулоны направляются на склад.
Склады ограничены по вместимости и максимальной допустимой ширине рулонов.
Потом со склада они отправляются на нарезку.
Роль склада в такой системе нельзя трактовать только как пассивное место хранения. Склад выступает буфером, который позволяет распределить по времени выпуск полуфабриката на БДМ и операции резки на ПРС/БРС. С одной стороны, наличие буфера увеличивает гибкость планирования, позволяя сначала накопить нужный набор тамбуров, а затем выполнить более удобную последовательность раскроев. С другой стороны, ограниченная вместимость и ограничения по ширине делают склад самостоятельным фактором выполнимости: формально хороший план по отходам может оказаться нереализуемым, если в некоторые моменты времени он требует хранения слишком большого числа рулонов или слишком широких тамбуров.

% \subsection{Критерии качества и конфликт целей}\label{subsec:theor_ch/criteria}
На практике критерии качества плана противоречивы: снижение отхода часто требует более «тонкой» подгонки раскроев, что может увеличивать число раскроев и количество переналадок.
В свою очередь, минимизация переналадок обычно достигается укрупнением партий и упрощением последовательности раскроев, что может приводить к росту отходов или увеличению времени ожидания заказов.
Поэтому модель должна позволять сравнивать планы по нескольким показателям и выбирать компромиссное решение.

Так как нужно учитывать большое количество ограничений, то используются планы производства.
Планы производства включают в себя планы производства бумаги на БДМ, планы нарезки на ПРС и БРС.
Есть критерии, которые определяют эффективность плана производства:
\begin{enumerate}
    \item минимизацию отходов, что позволяет снизить издержки и уменьшить воздействие на окружающую среду \cite{haessler1991};
    \item минимизацию числа раскроев, так как каждый раскрой связан с дополнительными затратами времени и материальных ресурсов\cite{cerqueira2021, vanderbeck2000};
    \item минимизацию перестановок ножей, отражающих трудоёмкость переналадки и простой оборудования \cite{cyberleninka_slitting_params, mobasher2013}.
\end{enumerate}
Оптимизация именно этих критериев позволяет найти баланс между экономической эффективностью, качеством исполнения заказов и соблюдением технологических ограничений.
Выбор именно этих критериев связан с практикой оперативного планирования на предприятии. Отход материала непосредственно влияет на себестоимость продукции и расход сырья, поэтому выступает главным экономическим показателем качества плана. Число раскроев характеризует организационную сложность исполнения плана: чем больше отдельных схем требуется выполнить, тем выше нагрузка на персонал и оборудование. Число перестановок ножей отражает скрытые потери времени, связанные с переналадкой, и потому выступает естественной мерой технологической устойчивости решения. Совместное рассмотрение этих критериев позволяет перейти от абстрактной минимизации отходов к оценке плана как производственно реализуемого объекта.

Особую сложность вызывает необходимость согласования планов работы БДМ и ПРС при наличии множества заказов, отличающихся по ширине, плотности, марке бумаги, срокам и дополнительным характеристикам \cite{keskinocak2002}.
В модели вводится правило упорядочивания марок бумаги в производстве, так как смена марки связана с технологическим простоем и расходом сырья.
Порядок марок может задаваться технологами до составления планов производства.
Для складов задаются ограничения по вместимости: они отражают реальные лимиты на хранение полуфабрикатов, обеспечивая, чтобы план не приводил к переполнению и сбоям в логистике.
Таким образом, планирование в рассматриваемой постановке носит интегрированный характер. Решение о том, как раскраивать тамбуры, нельзя принимать без учёта того, на каких БДМ они будут произведены, когда поступят на склад и в какой последовательности смогут быть переработаны на резательном оборудовании. Напротив, изолированное построение только плана раскроя или только календарного плана приводит к риску получить решение, которое оптимально по отдельному локальному критерию, но не выполнимо в общей технологической цепочке. Именно необходимость согласования стадий выпуска, хранения и резки определяет структуру последующей математической модели.

Таким образом, составление оптимального плана производства с учётом всех производственных параметров и ограничений является ключевым инструментом повышения эффективности работы предприятия.
Это обеспечивает не только рост конкурентоспособности компании, но и вклад в устойчивое развитие отрасли за счёт ресурсосбережения и снижения негативного экологического воздействия.
Математическая модель и алгоритм позволяют сформировать план производства и нарезки, минимизирующий отходы, переналадки и издержки, а также гарантировать своевременное выполнение заказов \cite{keskinocak2002}.

В математической модели нужно включить следующие факторы:
\begin{itemize}
    \item заказы представляют собой рулоны и имеют срок, ширину, количество рулонов, количество слоев, плотность, марку бумаги.
В зависимости от ширины они должны быть выпущены на ПРС или БРС;
    \item марки бумаги должны выставляться на БДМ в соответствии с заданным порядком.
Смена марки затрачивает время одинаковое для всех БДМ;
    \item для ПРС и БРС задается единый размер кромки тамбура.
Половина этого размера должна быть срезана слева, а другая половина справа.
При самой первой нарезке можно выставить нож вначале на отступ соответствующий половине размера кромки и не передвигать его при последующих нарезках.
А другой край будет срезаться при нарезки рулонов, но важно чтобы нож не заходил на эту половину размера кромки в конце рулона;
    \item время начала производства;
    \item БДМ, каждый из которых имеет ширину выпускаемого тамбура, скорость производства бумаги, время смены плотности;
    \item ПРС и БРС, каждый из которых имеет максимальную ширину нарезаемого рулона, количество ножей, скорость нарезки одного раскроя с учетом количества слоев, скорость перестановки ножей;
    \item склады рулонов для нарезки, которые имеют максимальную вместимость по количеству рулонов и  максимально допустимую ширину рулона;
    \item момент появления тамбура на складе соответствует окончанию его производства на БДМ;
    \item при выполнении нарезки на ПРС или БРС со склада изымаются необходимые рулоны;
    \item БДМ и ПРС могут работать параллельно, если их процессы не конфликтуют по состоянию склада.
    \item итоговый план производства, который учитывает все ограничения и отражает последовательность действий для выполнения всех заказов.
\end{itemize}

% \subsection{Принятые допущения и ограничения применимости}\label{subsec:theor_ch/assumptions}
При формализации задачи используются допущения, позволяющие построить модель, пригодную для алгоритмического поиска решений.
Принятые допущения не устраняют производственную специфику задачи, а переводят её в форму, пригодную для формального анализа и вычислительного эксперимента. Иными словами, они отделяют существенные факторы, которые необходимо учитывать для получения корректного плана, от второстепенных деталей, усложняющих модель без заметного выигрыша в качестве решения. Такой подход соответствует общепринятой логике построения прикладных оптимизационных моделей: сначала фиксируется управляемое ядро постановки, а затем оценивается, насколько полученная модель адекватна исходному производственному процессу.
\begin{itemize}
    \item Модель ориентирована на планирование по ширинам рулонов и учитывает кромку как технологический отход.
    \item Порядок марок бумаги в выпуске полуфабриката задаётся технологически и рассматривается как входной параметр планирования.
    \item Параметры оборудования (скорости производства и резки, время перестановки ножей, время смены марки/плотности) считаются известными и постоянными на горизонте планирования; возможные колебания учитываются на этапе экспериментальной валидации.
    \item Склад моделируется укрупнённо ограничениями на вместимость и допустимую ширину; детальные аспекты адресного хранения и внутренней логистики не рассматриваются.
    \item Модель нацелена на построение выполнимого плана для набора заказов; вопросы долгосрочного стратегического планирования и прогнозирования спроса находятся вне рамок рассмотрения.
    \item Горизонт планирования составляет один месяц, и формирование плана производства выполняется раз в месяц, что соответствует типичной практике оперативного планирования на предприятиях ЦБП.
\end{itemize}

В системе управления предприятием рассматриваемая задача относится именно к уровню оперативного планирования. В отличие от стратегических решений, связанных с выбором состава оборудования, ассортиментной политики или долгосрочных производственных мощностей, здесь предполагается, что парк БДМ, ПРС, БРС и складов уже задан. В отличие от укрупнённого календарно-производственного планирования, где часто оперируют суммарными объёмами выпуска по периодам, настоящая постановка детализирует последовательность выпуска полуфабриката, конкретные раскрои и порядок выполнения операций на оборудовании. Иными словами, задача решает вопрос не о том, что предприятие в целом должно произвести в течение длительного периода, а о том, как именно организовать исполнение уже известного набора заказов на заданном горизонте с учётом технологических ограничений.
Входными данными для такой постановки служат заказы клиентов, параметры БДМ, характеристики ПРС и БРС, ограничения складов, технологически заданный порядок марок бумаги, а также начальные условия планирования. Эти данные формируют исходное состояние системы и задают множество допустимых решений. Выходом же является не один агрегированный показатель, а согласованный набор трёх компонент: план производства бумаги $\overline{\overline{X}}$, план нарезки $\overline{\overline{Y}}$ и расписание $\overline{\overline{Z}}$. Только их совместное построение даёт результат, пригодный для практического использования, поскольку инженер или диспетчер должен получить не абстрактную оценку качества, а последовательность действий, интерпретируемую на уровне реального производственного процесса.
Выбор месячного горизонта планирования также имеет прикладное обоснование. С одной стороны, этого периода достаточно, чтобы учесть совокупность заказов, сгладить локальные колебания загрузки и согласовать работу нескольких участков производства. С другой стороны, слишком длинный горизонт снижал бы достоверность исходных данных: меняются приоритеты заказов, уточняются сроки, могут корректироваться производственные ограничения. Поэтому месячный горизонт представляет собой разумный компромисс между полнотой учёта производственной ситуации и устойчивостью исходной информации, что соответствует практике оперативного управления на предприятиях целлюлозно-бумажной промышленности.

% \subsection{Выводы по разделу}\label{subsec:theor_ch/problem_conclusion}
В разделе выделены элементы производственной цепочки (выпуск полуфабриката, резка, складирование) и сформулированы ключевые критерии качества плана.
Далее в обзоре литературы рассматриваются подходы, позволяющие учитывать данные особенности в постановках задач раскроя и календарного планирования, а затем вводится формальная модель, которая используется в последующих главах для построения алгоритма и программной реализации.
Тем самым раздел задаёт переход от содержательного описания производственной системы к формальной постановке исследовательской задачи. Выделенные технологические ограничения и критерии качества далее используются как основание для выбора математической модели, структуры алгоритма и программной реализации. Это обеспечивает логическую связность диссертации: обзор предметной области, формальная модель, численные методы и программный комплекс описывают одну и ту же задачу на разных уровнях детализации.


\section{Обзор литературы}\label{sec:theor_ch/literature_review}
% TODO: добавить всех-всех отсюда file:///home/vladislav/Downloads/autoref-matematicheskie-modeli-i-chislennye-metody-v-programmnom-komplekse-optimizatsii-planirovaniy.pdf
Цель данного обзора — показать развитие классической задачи линейного раскроя и связанных с ней постановок, а также выделить направления, наиболее близкие к задачам целлюлозно-бумажного производства: учёт потерь при продольной резке, ограничений по переналадке (перестановка ножей, смена схем), интеграция раскроя с календарным планированием и со складскими ограничениями.
Особое внимание уделяется постановкам, в которых помимо минимизации отходов рассматриваются эксплуатационные критерии (число раскроев/операций и переналадок), поскольку именно они определяют практическую реализуемость производственного плана.

% \subsection{Классическая задача линейного раскроя и её обобщения}\label{subsec:theor_ch/csp_classic}

Задача линейного раскроя (ЗЛР) относится к числу классических задач комбинаторной оптимизации и широко применяется при планировании производства в отраслях, где исходный материал поступает в виде заготовок фиксированной длины или ширины, а выходная продукция представлена наборами требуемых размеров.
К таким отраслям относятся, в частности, металлургия, деревообработка, текстильная промышленность и целлюлозно-бумажное производство.
Фундаментальные формальные постановки задачи линейного раскроя предложены П.~Гилмором и Р.~Гомори \cite{gilmore1961linear, gilmore1963linear}, где задача представлена в виде модели целочисленного линейного программирования, а также разработан подход генерации столбцов, позволяющий эффективно решать задачи большой размерности за счёт построения только необходимых раскроев; дальнейшее развитие и систематизация этого направления подробно изложены в \cite{desaulniers2005_column_generation, scheithauer2017_cutpack_optimization, goulimis1990_cutstock_optimal}.
Русскоязычные учебные и обзорные издания также фиксируют базовые модели ЗЛР и связь с линейным/целочисленным программированием, включая вопросы построения шаблонов и интерпретации двойственных оценок \cite{агишева2010линейное,зенков2017методы, dantzig1963_lp_extensions, bazaraa2011_lp_network, hillier2021_intro_or, ahuja1993_network_flows}.
Дальнейшее развитие исследований привело к формированию широкого класса задач раскроя и упаковки.
Для систематизации постановок и выбора методов решения существенную роль сыграла типология, предложенная Х.~Дикхоффом \cite{dyckhoff1990typology}, а также её развитие и уточнение в более поздних работах, в частности в \cite{wascher2007typology}; обзорные библиографические и типологические материалы по cutting and packing также представлены в \cite{dyckhoff1992_typology_biblio, sweeney1991_cutpack_review}.
Указанные классификации позволяют корректно отнести прикладные варианты раскроя, включающие дополнительные технологические ограничения, к соответствующим подклассам и тем самым обосновать применимость тех или иных методов решения.

С точки зрения вычислительной сложности задача раскроя относится к трудным комбинаторным задачам; на практике это означает необходимость компромисса между точностью и временем расчёта, особенно при росте числа заказов и расширении набора ограничений \cite{garey1979computers, papadimitriou1998_combinatorial_opt}.
Поэтому в реальных производственных постановках широкое распространение получили либо точные методы с использованием декомпозиции, либо эвристические и метаэвристические алгоритмы, обеспечивающие приемлемое качество решений при ограниченном времени.

В более широком контексте постановки задач раскроя, а также их интеграция с планированием производства и логистикой, естественным образом опираются на аппарат дискретной и комбинаторной оптимизации и связанные модели на графах \cite{струченков2017методы,Rajgorodsky2012, wolsey1998_integer_programming, schrijver1998_theory_lip}.



% \subsection{Точные методы решения: декомпозиция и метод ветвей и цен}\label{subsec:theor_ch/exact_methods}

Для задачи линейного раскроя и её обобщений важную роль играют точные методы, основанные на декомпозиции и последовательном уточнении решения.
Подходы, использующие декомпозицию Дантцига--Вульфа и генерацию столбцов, позволяют получать сильные нижние оценки и формировать решения в пространствах большой размерности \cite{desaulniers2005_column_generation, dantzig1963_lp_extensions}.
Практические аспекты построения алгоритмов типа «ветвей и цен» (ветвление в сочетании с генерацией столбцов), а также рекомендации по их эффективной реализации подробно обсуждаются в работе Ф.~Вандербека \cite{vanderbeck2000integer}.
Современный обзор математических моделей и точных алгоритмов для задач раскроя и родственных задач упаковки представлен в \cite{delorme2016bin}, где рассматриваются типовые целочисленные формулировки, схемы построения раскроев и методы улучшения вычислительной устойчивости; для двумерных постановок аналогичные идеи глобального поиска обсуждаются также в \cite{huang2022_global_2dcsp}.
Для подзадач генерации схем раскроя существенную роль играют рюкзачные постановки и связанные алгоритмы динамического программирования \cite{martello1990_knapsack}.
Обоснование применения аппарата целочисленного программирования к производственному планированию и смежным комбинаторным задачам также приведено в \cite{pochet2006_prodplan_mip, wolsey1998_integer_programming}.
В отечественной литературе акцентируется, что включение технологических ограничений (например, ограничения на ресурсы, смешанные целочисленные переменные, многопроизводственные особенности) приводит к усложнению моделей и обычно требует комбинации точных и приближённых процедур \cite{листопад2017некоторые,архипов2014математические,васильев2016методы}.

Точные методы особенно полезны для получения эталонных решений на задачах малой и средней размерности, а также для построения нижних оценок, позволяющих сравнивать качество эвристических решений.
Однако при добавлении технологических ограничений (переналадки, очередность выполнения схем, складирование, сроки) размерность модели и сложность подзадач генерации колонок существенно возрастают, что ограничивает применение точных алгоритмов в режиме оперативного планирования.

% \subsection{Приближённые методы: эвристики и метаэвристики}\label{subsec:theor_ch/heuristics}

В прикладных постановках, возникающих в промышленности, размерность задачи и число технологических ограничений зачастую делают применение точных методов ограниченным по вычислительным ресурсам.
В связи с этим широко используются приближённые методы, в том числе эвристики (конструктивные и улучшающие), а также метаэвристики, включая локальный поиск и гибридные алгоритмические схемы.
В работе Т.~Урбана \cite{urban2015methods} исследованы приближённые методы решения задачи раскроя на основе адаптивного поиска и показана эффективность подобного подхода для получения качественных решений при ограниченном времени расчёта.
Существенная часть современных прикладных решений основывается на комбинации конструктивной генерации схем раскроя и процедур улучшения решения.
Для задач упаковки и раскроя в условиях ограниченных ресурсов и жёстких технологических правил широко обсуждаются улучшенные жадные процедуры и локальные перестройки, применимые в промышленной практике \cite{чеканин2012оптимизация,панюкова2012оптимизация}.

В качестве универсального семейства метаэвристик часто используются генетические алгоритмы, сформировавшиеся как общий подход к оптимизации в сложных пространствах решений \cite{goldberg1989genetic, holland1992_adaptation, goldberg1989_ga_search}.
Для задач раскроя генетические алгоритмы привлекательны тем, что позволяют естественно кодировать раскрои и порядок их выполнения, а также легко дополняются механизмами учёта отраслевых ограничений (переналадки, складирование, многокритериальность).
Практическая эффективность генетического алгоритма для ЗЛР демонстрировалась, в частности, в \cite{haessler1991}, а дальнейшие исследования развивали идеи многокритериальной оптимизации и настройки операторов эволюции под специфику производственных процессов \cite{araujo2014genetic,golfeto2009genetic,lai1997developing}.
Базовые принципы кодирования решений и настройки операторов генетических алгоритмов систематизированы в учебных и обзорных работах \cite{gladkov2010,kureichik1998, eiben2015_intro_ec, gendreau2019_handbook_metaheur}, а развитие самоадаптивных схем эволюционного поиска и их применение к сложным задачам оптимизации обсуждается в \cite{semenkina2013,diveev2017, deb2001_moo_ea}.
В контексте многокритериальности и раскроя дополнительно отмечаются подходы, ориентированные на построение компромиссных решений для крупных экземпляров \cite{архипов2015применение}.

% \subsection{Потери материала при продольной резке и отраслевые постановки}\label{subsec:theor_ch/trim_loss}

Для целлюлозно-бумажной промышленности характерны постановки, ориентированные на минимизацию потерь материала при продольной резке рулонных заготовок, включая потери на обрезь и недоиспользование ширины.
Обзор соответствующих постановок и методов решения, а также обсуждение практических аспектов минимизации потерь приведены в работе А.~Хинксмана \cite{hinxman1980trimloss}.
В отечественной литературе систематическое изложение математических моделей и оптимизационных подходов для предприятий ЦБП представлено в монографиях А.\,В.~Воронина и В.\,А.~Кузнецова \cite{воронин2000математические,воронин2000прикладные}, где рассмотрены прикладные задачи планирования производства и управления материальными потоками, включая задачи раскроя и формирования производственных программ.
В диссертации Р.\,В.~Воронова \cite{воронов2004диссертация} разработаны математические модели и методы автоматизированного планирования производства бумаги, включающие распределение заказов между БДМ, формирование схем раскроя и построение объёмно-календарных планов с~учётом технологических ограничений.
Развитием данного направления является диссертация А.\,Р.~Урбана \cite{урбан2015диссертация}, в~которой предложены математические модели и~численные методы оптимизации планирования производства бумаги, реализованные в~программном комплексе; рассмотрены задачи компоновки нестандартных съёмов тамбуров и~учёт дополнительных технологических параметров.
В отличие от абстрактной постановки задачи линейного раскроя, отраслевые задачи минимизации потерь включают дополнительные компоненты: ограничения по настройке ножей, зависимость затрат от смены формата, наличие ограничений по складам и допустимой перепроизводству/недопроизводству, а также требования по срокам выполнения заказов.

Важным аспектом отраслевых задач является также порядок применения раскроев, поскольку изменение схемы ножей и форматные переходы приводят к простоям и дополнительным затратам.
Поэтому на практике целесообразно учитывать не только состав раскроев, но и структуру последовательности их выполнения.
Применительно к бумагоделательной промышленности А.\,Р.~Урбаном \cite{ромолдович2015методы} предложены методы решения задачи компоновки нестандартных съемов тамбуров, учитывающие специфику формирования планов продольной резки при наличии технологических ограничений.
Математические модели и алгоритмы решения задач оптимального раскроя полосы, применимые к рулонным материалам, исследованы в работе Р.\,В.~Сошкина~\cite{сошкин2009математические}.
Ч.~Чен и~соавт.~\cite{chen2019skiving} исследовали задачу склейки и~раскроя нестандартных рулонов (skiving and cutting stock problem) в~бумажной и~плёночной промышленности, где остатки от~предыдущих нарезок склеиваются продольно и~затем раскраиваются на~готовую продукцию; предложены модели ЦЛП и~эвристика последовательной коррекции значений.

\noindent Отдельное направление составляет учёт затрат на переналадку (setup costs) и зависимость целевой функции от числа смен схем раскроя и перестановок ножей.
Подобные постановки рассматриваются в работах, где расходы на переналадку включаются в оптимизационную модель наряду с отходом материала \cite{mobasher2013,martina2022}.
Для производственных линий ЦБК такие компоненты имеют принципиальное значение, поскольку переналадка влияет не только на трудоёмкость, но и на пропускную способность участка резки.

Практические технологические аспекты, связанные с параметрами режущего инструмента и требованиями к качеству обработки, также отражаются в прикладных исследованиях смежных процессов (например, высечки и резки рулонных материалов), что может использоваться при обосновании ограничений на число ножей и режимы переналадки \cite{bobrov2016}.

% \subsection{Интеграция раскроя с календарным планированием и очередностью выполнения работ}\label{subsec:theor_ch/integration_scheduling}

В производственных системах с несколькими взаимосвязанными стадиями (выпуск полуфабриката, резка, складирование и отгрузка) задача раскроя тесно связана с задачей формирования производственного расписания.
Учет временных ограничений и переналадок оборудования приводит к необходимости совместного рассмотрения минимизации потерь материала и минимизации затрат времени/переналадок.
Пример постановок календарного планирования операций резки на нескольких параллельных машинах представлен в \cite{giannelos2001parallel, pierce1964_paper_scheduling}.
В работе В.\,В.~Полякова и соавт.\ предложена модель оперативного планирования выпуска бумаги с использованием интервальных переменных, позволяющая учитывать неопределённость технологических параметров при составлении производственного расписания~\cite{поляков2005оперативное}.

Подходы, ориентированные на одновременное решение задач минимизации потерь и формирования расписания в целлюлозно-бумажной отрасли, изучались, в частности, в \cite{harjunkoski1996trimloss, westerlund1998trimloss_scheduling}, где представлены целочисленные формулировки и обсуждены вычислительные особенности их применения.
Отдельно рассматриваются постановки, в которых задача раскроя дополняется ограничениями на очередность выполнения схем раскроя; пример интегрированной постановки «раскрой + очередность» представлен в \cite{yanasse2007integrated}.

Для практического планирования существенную роль играют ограничения, которые связывают расписание резки с доступностью полуфабрикатов и состоянием склада, а также обеспечивают выполнимость планов при параллельной работе нескольких машин.
В результате в рамках интегрированных постановок формируется не только набор схем раскроя, но и порядок их исполнения, согласованный с выпуском и логистикой полуфабрикатов, что особенно важно для непрерывных производств.

% \subsection{Интеграция раскроя с планированием выпуска партий и многопериодные постановки}\label{subsec:theor_ch/integration_lotsizing}

При наличии связи между выпуском полуфабриката и последующей резкой, а также при ограничениях складирования возникает необходимость интеграции задачи раскроя с задачами планирования выпуска партий и управления запасами.
Такой класс постановок рассматривается в работах, посвящённых совместному планированию объёмов производства и раскроя, например в \cite{nonas2000combined, quadt2004_lotsizing_flexflow, katok2005_lotsizing_assembly}.
Современная классификация интегрированных постановок «планирование партий + раскрой», а также обзор литературы приведены в \cite{melega2018review}.
Вопросы планирования партий при случайном выходе продукции также анализируются в обзорной работе \cite{yano1989_random_yields_review}, что важно для понимания границ применимости детерминированных моделей производства.
Для целлюлозно-бумажной промышленности известны прикладные модели, в которых целевая функция учитывает одновременно затраты производства/хранения и потери при раскрое \cite{poltroniere2016tema}.
В~работе А.~Леан и Ф.~Толедо \cite{leao2016decomposition} исследованы методы декомпозиции Дантцига–Вульфа для интегрированной задачи раскроя и~планирования партий в~бумажной промышленности с~учётом параллельных машин, времени переналадки и~управления запасами; показано, что генерация столбцов без декомпозиции по периодам даёт лучшие допустимые решения при~низких вычислительных затратах.
Для интеграции раскроя, многопериодного планирования и последовательностезависимых переналадок представляет интерес также диссертационная работа \cite{ditri1994_cutstock_scheduling_thesis}.

Интеграция с управлением запасами позволяет формализовать компромисс между немедленным выполнением резки и накоплением полуфабриката для повышения качества раскроя (снижение отхода и числа переналадок).
В таких моделях возникают дополнительные ограничения по ёмкости складов, динамике запасов и синхронизации стадий производства, что напрямую соответствует реальным процессам на целлюлозно-бумажных предприятиях.

% --- начало расширенного раздела 1.1 ---
% \paragraph{Колонковая генерация и линейное программирование.}
Классическим методом решения задачи линейного раскроя является линейно‑программная модель с генерацией колонок, предложенная Гилмором и Гомори в 1960‑х гг.
\cite{gilmore1961linear,gilmore1963linear}.
Алгоритм строит мастер‑проблему, переменные которой соответствуют допустимым раскроям, и последовательно генерирует новые колонки (раскрои) с помощью задачи на рюкзак (динамического программирования) \cite{струченков2017динамическое,лааксонен2019олимпиадное}.
В дальнейшем подход был расширён: Вандербек предложил неявную генерацию колонок и комбинирование её с процедурой ветвления \cite{vanderbeck2000integer}; Делорме и соавт.\ провели сравнительный анализ эффективности методов \emph{branch-and-price} для задач раскроя и упаковки \cite{delorme2016bin}, а Янассе и Пинто изучили интеграцию раскроя с планированием работы оборудования \cite{yanasse2007integrated}.
Эти работы показали, что колонковые методы хорошо работают на задачах средней размерности и пригодны для построения практически оптимальных решений.

% \paragraph{Многокритериальные модели.}
Во многих практических ситуациях необходимо минимизировать одновременно несколько критериев (например, отходы бумаги, количество переналадок на станках, число раскроев).
Бансал применил эволюционный алгоритм для одновременной минимизации отходов и числа раскроев, продемонстрировав на примерах, что многокритериальные алгоритмы позволяют находить компромиссные решения \cite{araujo2014genetic}.
Силва и соавт.\ предложили генетический алгоритм для одноразмерной задачи линейного раскроя (ЗЛР) с двумя критериями – потерей материала и числом операций \cite{golfeto2009genetic}.
Токсари и Байкасоглу разработали алгоритм имитации отжига, который учитывает и отходы, и время работы оборудования \cite{lai1997developing}.
Суджианто и~Ян предложили метаэвристику симбиотического поиска (SOS) для двухкритериальной задачи раскроя арматуры с~одновременной минимизацией отходов и~числа схем раскроя, подтвердив эффективность метаэвристического подхода к~многокритериальным постановкам \cite{sugianto2022rebar}.
Такие исследования подчёркивают, что многокритериальные постановки более полно отражают реальные потребности производства ЦБК, где важно не только сократить отход, но и снизить количество переналадок и время выполнения заказов.

% \paragraph{Эволюционные и эвристические подходы.}
Помимо линейно‑программных методов, значительное внимание уделяется эвристикам и метаэвристикам.
Ещё в 1990‑х гг.
Хесслер и Берда показали эффективность генетического алгоритма для задачи линейного раскроя (ЗЛР) \cite{haessler1991}.
В последующих работах генетические алгоритмы успешно применялись к более сложным постановкам, включая задачи с изменяемым набором оборудования и расписанием \cite{kureichik1998,genetics_2,evolution_1,semenkina2013}.
Обзор Мелега и соавт.\ подчёркивает, что эволюционные методы особенно полезны при большом числе технологических ограничений, когда методы линейного программирования становятся непрактичными \cite{melega2018review}.
Таким образом, выбор метода зависит от размерности и сложности задачи: для небольших примеров эффективнее точные линейные методы, а при масштабных ограничениях – генетические алгоритмы, что будет подробно рассмотрено в последующих главах.
% --- конец расширенного раздела 1.1 ---


% \subsection{Выводы по обзору и место настоящей работы}\label{subsec:theor_ch/lit_conclusion}

Таким образом, в литературе подробно изучены: (i) классическая задача линейного раскроя и методы её решения, (ii) отраслевые постановки минимизации потерь при продольной резке рулонных материалов, (iii) постановки, учитывающие очередность выполнения схем раскроя и переналадки, (iv) интеграция раскроя с задачами календарного планирования и планирования выпуска партий.

Вместе с тем для задач целлюлозно-бумажного производства сохраняется разрыв между моделями, ориентированными преимущественно на минимизацию потерь материала, и моделями, учитывающими согласованное формирование производственного расписания для взаимосвязанных стадий (выпуск полуфабриката, резка, складирование).
В~работах Р.\,В.~Воронова \cite{воронов2004диссертация} и А.\,Р.~Урбана \cite{урбан2015диссертация} заложены основы автоматизации планирования в~ЦБП, включая формирование схем раскроя, учёт переналадок ножей и~многокритериальную оптимизацию.
Вместе с~тем в~указанных работах не~рассматривались одновременный учёт нескольких ПРС и~БРС, многослойные заказы, динамика складских запасов как ограничение при формировании расписания, а~также построение согласованного расписания выпуска и~нарезки на~основе модели расписания.
Настоящая работа направлена на сокращение указанного разрыва за счёт разработки модели и алгоритмического подхода, обеспечивающих согласование решений по раскрою и формированию расписания с учётом технологических и временных ограничений.
Следовательно, новизна рассматриваемой постановки определяется не изолированным добавлением одного нового ограничения, а объединением нескольких факторов в рамках единого глобального плана. Существенным является именно совместный учёт разнородных аспектов: нескольких типов резательного оборудования, последовательности марок, многослойности заказов, динамики складов и календарной согласованности операций. Такая комбинация факторов изменяет не только вид ограничений, но и сам характер решения: вместо статического набора раскроев требуется построить связанный производственный сценарий, в котором каждая операция согласована с остальными по ресурсам и времени.

Сравнение элементов постановки в~диссертационных работах по~данной тематике приведено в~таблице~\ref{tab:lit_comparison}.

\begin{small}
\begin{longtable}{|p{0.22\linewidth}|c|c|c|c|c|c|}
\caption{Сравнение элементов постановки в~работах по~планированию в~ЦБП}%
\label{tab:lit_comparison}\\
\hline
\textbf{Элемент постановки} &
\textbf{\shortstack{Воро-\\нов\\(2004)}} &
\textbf{\shortstack{Ур-\\бан\\(2015)}} &
\textbf{\shortstack{Леан,\\Толедо\\(2016)}} &
\textbf{\shortstack{Чен\\и~др.\\(2019)}} &
\textbf{\shortstack{Суджи-\\анто,\\Ян (2022)}} &
\textbf{\shortstack{Наст.\\работа}} \\
\hline
\endfirsthead
\caption[]{Сравнение элементов постановки (продолжение)}\\
\hline
\textbf{Элемент постановки} &
\textbf{\shortstack{Воро-\\нов\\(2004)}} &
\textbf{\shortstack{Ур-\\бан\\(2015)}} &
\textbf{\shortstack{Леан,\\Толедо\\(2016)}} &
\textbf{\shortstack{Чен\\и~др.\\(2019)}} &
\textbf{\shortstack{Суджи-\\анто,\\Ян (2022)}} &
\textbf{\shortstack{Наст.\\работа}} \\
\hline
\endhead
\hline
\endfoot
Распределение заказов между БДМ & $+$ & $+$ & $+$ & $-$ & $-$ & $\boldsymbol{+}$ \\
\hline
Формирование схем раскроя & $+$ & $+$ & $+$ & $+$ & $+$ & $\boldsymbol{+}$ \\
\hline
Учёт переналадок ножей & $+$ & $+$ & $+$ & $-$ & $-$ & $\boldsymbol{+}$ \\
\hline
Многокритериальная ЦФ & $+$ & $+$ & $-$ & $-$ & $+$ & $\boldsymbol{+}$ \\
\hline
Календарное планирование & $+$ & $+$ & $+$ & $-$ & $-$ & $\boldsymbol{+}$ \\
\hline
Учёт сроков заказов & $+$ & $+$ & $+$ & $-$ & $-$ & $\boldsymbol{+}$ \\
\hline
Учёт марок и плотности & $+$ & $+$ & $+$ & $-$ & $-$ & $\boldsymbol{+}$ \\
\hline
Программный комплекс & $+$ & $+$ & $-$ & $-$ & $-$ & $\boldsymbol{+}$ \\
\hline
Генерация столбцов / ЛП & $+$ & $+$ & $+$ & $+$ & $-$ & $-$ \\
\hline
Метаэвристика & $-$ & $-$ & $-$ & $-$ & $+$ & $\boldsymbol{+}$ \\
\hline
Оптимизация порядка раскроев & $-$ & $+$ & $-$ & $-$ & $-$ & $\boldsymbol{+}$ \\
\hline
Нестандартные (склеенные) съёмы & $-$ & $+$ & $-$ & $+$ & $-$ & $-$ \\
\hline
Несколько единиц оборудования & $-$ & $+$ & $+$ & $-$ & $-$ & $\boldsymbol{+}$ \\
\hline
Интеграция раскроя с~планированием партий & $-$ & $-$ & $+$ & $-$ & $-$ & $-$ \\
\hline
Управление запасами / складами & $-$ & $-$ & $+$ & $-$ & $-$ & $\boldsymbol{+}$ \\
\hline
Несколько ПРС и~БРС & $-$ & $-$ & $-$ & $-$ & $-$ & $\boldsymbol{+}$ \\
\hline
Многослойные заказы & $-$ & $-$ & $-$ & $-$ & $-$ & $\boldsymbol{+}$ \\
\hline
\textbf{Динамика складов} & $-$ & $-$ & $-$ & $-$ & $-$ & $\boldsymbol{+}$ \\
\hline
\textbf{Согласованное расписание} & $-$ & $-$ & $-$ & $-$ & $-$ & $\boldsymbol{+}$ \\
\hline
\textbf{Модель расписания} & $-$ & $-$ & $-$ & $-$ & $-$ & $\boldsymbol{+}$ \\
\hline
\textbf{ГА с~глобальной мутацией} & $-$ & $-$ & $-$ & $-$ & $-$ & $\boldsymbol{+}$ \\
\hline
\end{longtable}
\end{small}

Как видно из таблицы, предшествующие работы содержат многие существенные элементы планирования ЦБП: распределение заказов, формирование раскроев, учёт переналадок и~многокритериальность.
Настоящая работа расширяет эти результаты за~счёт одновременного учёта нескольких ПРС и~БРС, многослойных заказов, динамики складских запасов, а~также за~счёт построения согласованного расписания на~основе модели расписания, связывающей выпуск, складирование и~нарезку в~едином плане.
Практическое значение такого расширения состоит в том, что предприятие получает возможность оценивать не только качество отдельной схемы раскроя, но и выполнимость всего комплекса производственных решений. Если в классических постановках основное внимание сосредоточено на снижении отхода, то в предлагаемом подходе дополнительно учитывается, сможет ли выбранная схема быть реализована на доступном наборе машин, без переполнения складов и без нарушения сроков. Это особенно важно для использования модели в составе программного комплекса, где результат должен быть интерпретируем не только как оптимизационное решение, но и как основа для практического производственного плана.
Дополнительного пояснения требует именно необходимость интегрированной постановки. Если оптимизировать только раскрой, то можно получить схему с малыми отходами, которая окажется не обеспеченной полуфабрикатом нужной ширины, марки или плотности в требуемый момент времени. Если же рассматривать только расписание выпуска и работы оборудования без явной модели раскроя, то остаётся неясным, какие именно сочетания форматов будут выполняться на ПРС и БРС и насколько рационально будет использован материал. В обоих случаях решение оказывается частичным и не даёт полной картины реального производственного процесса.
Существенную роль играет и наличие складов, а также второго уровня резки на БРС. Эти элементы делают задачу не просто расширением классической задачи линейного раскроя, а качественно иной постановкой, в которой необходимо согласовывать материальные потоки, временные зависимости и распределение операций между разными типами оборудования. Склад вводит ограничения на промежуточное хранение и тем самым влияет на допустимость расписания, а БРС требует учитывать преобразование промежуточных рулонов в готовую продукцию через дополнительный этап обработки. В результате локально эффективные решения для отдельных участков не обязательно совместимы между собой.
Поэтому естественным объектом оптимизации становится глобальный план $G=(\overline{\overline{X}},\overline{\overline{Y}},\overline{\overline{Z}})$, объединяющий выпуск полуфабриката, раскрой и календарное согласование операций. Только такая форма представления позволяет одновременно контролировать материальную обеспеченность, технологическую допустимость и соблюдение сроков. В следующей главе эта идея получает строгое математическое выражение в виде системы переменных, целевой функции и ограничений, описывающих все три компоненты плана как части единой модели.

Отдельно отметим, что в отечественных работах автора и соавторов рассматривались генетические и комбинированные подходы к одноразмерным задачам раскроя с учётом практических критериев качества.
В частности, предложены алгоритмы генерации схем раскроя и подходы к настройке эвристик \cite{klimenko2024moscow}, исследованы генетические алгоритмы для постановок с допусками на объёмы выпуска \cite{klimenko2024_prin}, а также рассмотрены методы, ориентированные на одновременное снижение отходов и числа перестановок ножей \cite{klimenko2025ivdon,klimenko2025fruct}.
Представленная в диссертации постановка развивает данные направления за счёт явного согласования стадии выпуска полуфабриката, операций резки и ограничений складирования в рамках единого плана.

Дополнительно в отечественной литературе представлены прикладные исследования и обзоры по смежным вариантам задач раскроя и упаковки, включая постановки одномерного раскроя и практические методы оптимизации в производственных условиях \cite{васин2012решение,агапов2021оптимизация,емельянов1996раскрой,тимофеева2013генетический}.
Вопросам построения многофункциональных программных систем для автоматизации решения задач раскроя материалов и комплектования изделий посвящена работа В.\,А.~Кузнецова и соавт.~\cite{кузнецов2008многофункциональная}.

Настоящая работа предлагает модель, в которой одновременно учитываются: требования по срокам, ограничения по складам, по переналадкам ножей и очередности схем раскроя – при одновременной минимизации нескольких критерия эффективности.
В известных нам публикациях такая комплексная постановка ранее не рассматривалась.
Тем самым обзор литературы обосновывает выбор направления исследования и показывает место настоящей работы среди существующих подходов. С одной стороны, работа опирается на развитый аппарат классической задачи раскроя, методов целочисленного программирования и метаэвристик. С другой стороны, предлагается постановка, ориентированная именно на интеграцию раскроя, выпуска полуфабриката, складирования и календарного планирования в условиях целлюлозно-бумажного производства. Это создаёт основу для последующего изложения математической модели в главе~\ref{ch:model} и численных методов её решения в главе~\ref{ch:algorithm}.

\FloatBarrier
